\section{ثالثًا: التصورات النظرية لأسلوب المعاملة الوالدية المتسلط}
يُعدّ أسلوب المعاملة الوالدية من أكثر المتغيّرات النفسية والاجتماعية تأثيرًا في تشكيل شخصية الطفل
وتطوّر مفهومه عن ذاته خلال مرحلة المراهقة. وتقتصر هذه الدراسة على استكشاف إدراك المراهق للأسلوب
المتسلط تحديدًا، وهو ما يستدعي التمييز الدقيق بين مصطلحين قد يلتبسان في الترجمة العربية: الأسلوب
المتسلط \en{(Authoritarian)} -- وهو محور هذا البحث -- القائم على ضبطٍ صارم مصحوبٍ بانخفاضٍ في
التقبّل العاطفي؛ والأسلوب السلطوي أو الديمقراطي التشاركي \en{(Authoritative)} القائم على ضبطٍ
متوازن مصحوبٍ بدفءٍ عاطفي مرتفع، وهو نمطٌ مختلف تمامًا في أثره رغم تشابه جذرهما اللغوي.

\subsection{تصوّر ديانا بومريند (Baumrind)}
تُعدّ ديانا بومريند من أبرز الباحثين في علم النفس النمائي المتخصّصين في سيكولوجية التنشئة الوالدية،
وهي الرائدة في تصنيف أساليب التنشئة، إذ ميّزت بين ثلاثة أنماط أساسية: الديمقراطي التشاركي
\en{(Authoritative)}، والمتساهل \en{(Permissive)}، والمتسلط \en{(Authoritarian)}. وتُعرّف الأسلوب
المتسلط بأنّه نمطٌ يميل إلى التحكّم في تشكيل سلوك الطفل وضبطه وتقييمه وفق معايير ثابتة وصارمة
يعتبرها الوالدان مطلقةً وغير قابلة للتفاوض، مع التأكيد على الطاعة وتقييد الاستقلالية واستخدام العقوبة،
وغيابٍ شبه تامّ للتفسير والحوار \en{(Baumrind, 1966)}. وأضافت لاحقًا أنّ هذا الأسلوب يتّسم بمستوى
عالٍ من الضبط والمطالب السلوكية يقابله مستوى منخفض من الاستجابة الوجدانية والدفء العاطفي
\en{(Baumrind, 1991)}.

\subsection{تصوّر ماكوبي ومارتن (Maccoby \& Martin)}
طوّر ماكوبي ومارتن نموذج بومريند بإضافة نمطٍ رابع هو الأسلوب المُهمِل \en{(Neglectful)}، وأعادا
تعريف الأسلوب المتسلط بوصفه نمطًا يتّسم بمستوى مرتفع من الضبط والتحكّم ومستوى منخفض من
الاستجابة الوجدانية، مؤكّدَين أنّ هذا الضبط الصارم في غياب الدفء العاطفي يُفضي إلى توتّراتٍ نفسية
لدى الأبناء، لا سيّما في مرحلة المراهقة \en{(Maccoby \& Martin, 1983)}. ويُضيف هذا التعريف على
بومريند بُعدًا تفسيريًا جوهريًا: فبينما وصفت بومريند الأسلوب المتسلط هيكليًا، ركّز ماكوبي ومارتن على
نتيجته النفسية المباشرة (التوتّر)، ممّا يمهّد لتساؤل هذا البحث عن انعكاسه على تقدير الذات.

\subsection{تصوّر باربر (Barber)}
يُميّز باربر بين نوعين من التحكّم الوالدي داخل الأسلوب المتسلط: \textbf{التحكّم السلوكي}
\en{(Behavioral Control)}، ويعني مراقبة سلوكيات الأبناء الخارجية وضبط أوقاتهم وعلاقاتهم؛
و\textbf{التحكّم النفسي} \en{(Psychological Control)}، ويشير إلى التحكّم في أفكار المراهق ومشاعره
وهويته الداخلية من خلال إثارة الذنب أو سحب الحبّ، وهو الأكثر ضررًا على الصحّة النفسية لأنّه يتدخّل
مباشرةً في بناء استقلالية المراهق الذاتية \en{(Barber, 1996)}.

وبالمقارنة بين التصوّرات السابقة يتّضح أنّها تمثّل بناءً نظريًا متكاملًا لفهم أسلوب المعاملة الوالدية
المتسلط؛ إذ قدّمت بومريند الإطار العام لهذا الأسلوب، ثمّ وسّع ماكوبي ومارتن هذا التصوّر بإبراز بُعدَي
الضبط والاستجابة الوالدية، في حين ركّز باربر على آليات التحكّم مميّزًا بين التحكّم السلوكي والنفسي.
ويكشف هذا التكامل أنّ الأسلوب المتسلط لا يقتصر على فرض الضبط، وإنّما يتضمّن أيضًا أنماطًا من
التفاعل قد تحدّ من استقلالية المراهق وتؤثّر في تفسيره لمعنى العلاقة الوالدية، وهو ما يبرّر اعتماد البحث
على المراجع الثلاثة معًا.

\subsection{أبعاد إدراك أسلوب المعاملة الوالدية المتسلط}
استنادًا إلى هذه التعريفات النظرية الثلاثة مجتمعةً، حُدّدت الأبعاد الفرعية الآتية التي شكّلت محاور شبكة
المقابلة:
\begin{itemize}[leftmargin=1.6em,itemsep=1pt]
\item \textbf{البُعد الأول -- الضبط والصرامة:} فرض قواعد صارمة وثابتة وتوقّع الطاعة الفورية دون مرونة
أو نقاش \en{(Baumrind, 1966)}.
\item \textbf{البُعد الثاني -- العقاب:} أسلوب الوالدين في الردع جسديًا أو نفسيًا/معنويًا
\en{(Baumrind, 1966; Barber, 1996)}.
\item \textbf{البُعد الثالث -- غياب الحوار والتفسير:} طبيعة الاتّصال الأحادي بين المراهق ووالديه دون
تفسير أو نقاش \en{(Baumrind, 1966)}.
\item \textbf{البُعد الرابع -- البرود العاطفي وقلّة التقبّل:} انخفاض مستوى الاستجابة الوجدانية والدفء
العاطفي \en{(Baumrind, 1966; Maccoby \& Martin, 1983)}.
\item \textbf{البُعد الخامس -- الرقابة والتدخّل:} درجة مراقبة الوالدين لتفاصيل حياة المراهق، ويقابل
«التحكّم السلوكي» عند باربر \en{(Barber, 1996)}.
\item \textbf{البُعد السادس -- كبت الاستقلالية والرأي:} تقييد قدرة المراهق على التعبير عن رأيه، ويتقاطع
مع «التحكّم النفسي» عند باربر \en{(Barber, 1996)}.
\end{itemize}

\subsection{التعريف الإجرائي لإدراك أسلوب المعاملة الوالدية المتسلط}
يُقصد بإدراك أسلوب المعاملة الوالدية المتسلط في هذه الدراسة الكيفيةُ التي يفسّر بها المراهق ممارسات
والديه ويمنحها معنى في ضوء خبراته الأسرية، وذلك بالاستناد إلى ستّة أبعاد هي: الضبط والصرامة،
والعقاب، وغياب الحوار والتفسير، والبرود العاطفي وقلّة التقبّل، والرقابة والتدخّل، وكبت الاستقلالية
والرأي.

\section{رابعًا: المفهوم الثاني -- تقدير الذات لدى المراهقين}
\section*{تمهيد}
يُعدّ تقدير الذات أحد المفاهيم الجوهرية في علم النفس، إذ يرتبط ارتباطًا وثيقًا بالصحّة النفسية للفرد
وتوافقه الشخصي والاجتماعي. ويزداد بناؤه حساسيةً في مرحلة المراهقة المتوسّطة تحديدًا نتيجة التفاعل بين
العوامل النمائية والعلاقة مع الوالدين؛ فما يُدركه المراهق من معاملة والديه المتسلطة هو المادة الخام التي
تتشكّل منها درجة تقديره لذاته، وفق ما تؤكّده الأدبيات النفسية.

\subsection{التناول النظري لمفهوم تقدير الذات}
\noindent\textbf{تعريف روزنبرغ (Rosenberg).}\ يُعرّف روزنبرغ تقدير الذات بأنّه الاتّجاه الإيجابي أو
السلبي الذي يحمله الفرد نحو ذاته كموضوعٍ كلّي، مؤكّدًا أنّ تقدير الذات المرتفع لا يعني أن يرى الفرد
نفسه أفضل من الآخرين بل يعني رضاه عن ذاته، في حين يعكس تقدير الذات المنخفض رفض الفرد لذاته
وعدم رضاه عنها \en{(Rosenberg, 1965)}؛ فهو يقدّم التعريف الأعمّ لتقدير الذات كاتّجاهٍ كلّي تنطلق منه
التعريفات اللاحقة.

\noindent\textbf{تعريف كوبر سميث (Coopersmith).}\ يلتقي كوبر سميث مع روزنبرغ في الفكرة العامّة
لكنّه يُضيف بُعدًا تركيبيًا أدقّ، إذ يُعرّف تقدير الذات بأنّه التقييم المعتاد والمستمرّ الذي يُجريه الفرد نحو
ذاته ويُعبّر عن الحدود التي يعتقد عندها أنّه كفءٌ وذو أهمّية وناجح وذو قيمة \en{(Coopersmith, 1967)}؛
فهو لا يكتفي بوصف التقييم كاتّجاهٍ عامّ بل يُفكّكه إلى مكوّناته، وهو ما يجعل تعريفه أقرب إلى منطق
الأبعاد الذي يحتاجه هذا البحث.

\noindent\textbf{تصوّر ماسلو (Maslow).}\ يضع ماسلو تقدير الذات ضمن هرم الحاجات الإنسانية مميّزًا
بين مستويين: تقدير الذات الداخلي (الكفاءة والإنجاز والاستقلالية) وتقدير الذات الخارجي (الاحترام
والاعتراف من الآخرين)، مؤكّدًا أنّ إشباع هذه الحاجة شرطٌ أساسي لتحقيق الذات \en{(Maslow, 1954)}؛
فغياب الاحترام في إطار المعاملة الوالدية المتسلطة (البرود العاطفي والعقاب) يُخلّ بإشباع الحاجة إلى
تقدير الذات الخارجي، ممّا ينعكس سلبًا على الشقّ الداخلي (الثقة بالنفس).

\noindent\textbf{تصوّر هارتر (Harter) لتقدير الذات لدى المراهق.}\ طوّرت هارتر هذا الطرح خصّيصًا في
سياق مرحلة المراهقة، مميّزةً بين عدّة مجالات يتشكّل منها تقدير الذات: المظهر الجسمي، والقبول
الاجتماعي من الأقران، والكفاءة الأكاديمية والأخلاقية، مؤكّدةً أنّ التقدير العامّ للذات هو محصّلة إدراك
المراهق لذاته في هذه المجالات مجتمعةً \en{(Harter, 1988)}. وبخلاف روزنبرغ وماسلو اللذين قدّما تصوّرًا
عامًّا لكلّ الأعمار، قدّمت هارتر تصوّرًا مُخصّصًا لمرحلة المراهقة، وهو ما يجعل تصنيفها الأقرب إلى فئة
هذا البحث العمرية.

ويتّضح من تحليل هذه التصوّرات أنّ تناول مفهوم تقدير الذات يتدرّج من العامّ إلى الأكثر تخصّصًا؛ إذ
ينظر إليه روزنبرغ اتّجاهًا عامًّا نحو الذات، بينما يفصّله كوبر سميث إلى مكوّنات، ثمّ يضعه ماسلو في إطار
الحاجات الإنسانية، في حين تقدّم هارتر تصوّرًا يفسّره في المراهقة عبر مجالاته المختلفة، وهو ما يبرّر
اعتماد الدراسة على أبعاد كوبر سميث وهارتر بوصفها الأقرب إلى طبيعة البحث وأهدافه.

\subsection{أبعاد تقدير الذات لدى المراهقين (الأبعاد المزدوجة)}
بما أنّ هذا البحث يسعى إلى استكشاف تقدير الذات من خلال إدراك المراهق لتسلّط والديه، فإنّ كلّ بُعدٍ
يُستكشَف بصيغةٍ مزدوجة تربطه بالبُعد المقابل له من أبعاد التسلط:
\begin{itemize}[leftmargin=1.6em,itemsep=1pt]
\item \textbf{تقدير الذات الشخصي/العامّ:} التصوّر الكلّي للمراهق نحو ذاته \en{(Rosenberg, 1965)}، في
علاقته ببُعدَي غياب الحوار وكبت الاستقلالية.
\item \textbf{تقدير الذات الأسري:} شعور المراهق بقيمته ومكانته داخل أسرته \en{(Coopersmith, 1967)}،
في علاقته ببُعدَي البرود العاطفي والعقاب.
\item \textbf{تقدير الذات الاجتماعي:} تصوّر المراهق لذاته في علاقاته مع أقرانه \en{(Harter, 1988)}، في
علاقته ببُعد الرقابة والتدخّل.
\item \textbf{تقدير الذات الأكاديمي:} شعور المراهق بكفاءته الدراسية \en{(Coopersmith, 1967; Harter,
1988)}، في علاقته ببُعد الضبط والصرامة.
\item \textbf{تقدير الذات الجسمي:} رضا المراهق عن مظهره الجسدي \en{(Harter, 1988)}، في علاقته ببُعد
البرود العاطفي وغياب التقبّل.
\end{itemize}

\subsection{التعريف الإجرائي لتقدير الذات}
يُعرَّف تقدير الذات إجرائيًا في هذه الدراسة بأنّه التقييم الذي يبنيه المراهق لذاته في مختلف جوانب حياته،
كما يظهر من خلال إدراكه لأسلوب المعاملة الوالدية المتسلط وتفسيره لخبراته الأسرية، ويُستدلّ عليه من
مضامين المقابلات نصف الموجّهة عبر خمسة أبعاد تحليلية هي: تقدير الذات الشخصي، والأسري،
والاجتماعي، والأكاديمي، والجسمي.

\section*{خلاصة الباب الأول}
بيّنت المراجعة النظرية أنّ مرحلة المراهقة، ولا سيّما الفئة العمرية (15--17 سنة)، تمثّل مرحلةً نمائية
تتّسم بتزايد قدرة المراهق على إدراك أساليب معاملة والديه، وبحساسيةٍ أكبر تجاه تأثيرها في تكوين صورته
عن ذاته. كما أظهرت الأدبيات أنّ الأسلوب الوالدي المتسلط يقوم على الضبط الصارم والعقاب وضعف
الحوار وانخفاض الدفء العاطفي والرقابة المفرطة وكبت الاستقلالية، وهي الأبعاد التي اعتمدتها الدراسة في
بناء شبكة المقابلة. كما أوضحت التصوّرات النظرية لتقدير الذات أنّه بناءٌ نفسي متعدّد الأبعاد يتشكّل من
خبرات الفرد وتفاعلاته الأسرية والاجتماعية، ممّا يبرّر اعتماد أبعاده الخمسة بوصفها الأنسب لفهم تجربة
المراهق. واستنادًا إلى هذا البناء النظري، ينتقل الباب الموالي إلى عرض الإطار المنهجي للدراسة: منهجها،
وعيّنتها، وأداة جمع معطياتها، وإجراءات تحليلها، بما ينسجم مع طبيعة الدراسة الكيفية وأهدافها.
